% ==============================================================================
% Section 1: Institutional Header, Abstract & Introduction
% ==============================================================================

\begin{center}
    % Pondicherry University Greyscale Crest
    \includegraphics[width=2.4cm]{figures/pondicherry_university_logo_grayscale.png}\\[8pt]

    {\large\textbf{DEPARTMENT OF STATISTICS \quad$\vert$\quad RAMANUJAN SCHOOL OF MATHEMATICAL SCIENCES}}\\[2pt]
    {\large\textbf{PONDICHERRY UNIVERSITY (A CENTRAL UNIVERSITY)}}\\[2pt]
    {\small Puducherry -- 605 014, India}\\[10pt]

    {\Large\bfseries DOWNSIDE RISK QUANTIFICATION AND DERIVATIVE HEDGING EFFICIENCY}\\[4pt]
    {\large\bfseries An Empirical Application of Value at Risk (VaR) and the Black--Scholes Model on ITC Limited}\\[8pt]

    {\small\textbf{Course:} MSQF 533 -- Financial Risk Management (Semester III) \quad$\vert$\quad \textbf{Academic Year:} 2026--2027}\\[2pt]
    {\small\textbf{Author:} N. Rohit Vedhanandh \quad$\vert$\quad \textbf{Reg. No.:} 25MSQUFPY0002 \quad$\vert$\quad \textbf{Date:} November 2026}\\[6pt]
    \rule{\linewidth}{0.8pt}
\end{center}

\begin{abstract}
\noindent
This study establishes a unified quantitative risk framework that links downside market risk quantification via Value at Risk ($\VaR$) and Expected Shortfall ($\CVaR$) directly to derivative hedging under the Black--Scholes--Merton option pricing model. Calibrating the framework to 246 daily trading sessions of \textbf{ITC Limited (NSE: ITC)} over a trailing 1-year sample period (November 2025 to November 2026) for 1 standard NSE F\&O market lot ($1,725$ shares at spot $S_0 = \text{\rupee}285.00$, total capital commitment $V_0 = \text{\rupee}4,91,625.00$), we reveal three primary empirical insights: First, daily log returns exhibit severe leptokurtosis ($\hat{K}_{\text{excess}} = +9.82$) and negative skewness ($\hat{S} = -0.92$), leading to overwhelming rejection of normality ($JB = 979.34, p < 10^{-10}$). Second, Gaussian models severely underestimate extreme tail risk: the 10-day $99\%$ Historical Simulation $\VaR$ reaches \textbf{\rupee 61,487.02} ($12.51\%$ of capital), producing an \textbf{11.70\% fat-tail wedge} over Parametric $\VaR$ (\rupee 55,048.61), while both pass Kupiec POF backtests within the Basel Green Zone. Third, neglecting ITC's $5.09\%$ trailing dividend yield introduces up to a $+9.76\%$ pricing error in standard Black--Scholes call options. Deploying an \textbf{ATM Protective Put} eliminates $75.14\%$ of extreme 10-day tail risk, while a \textbf{Zero-Cost Collar} delivers a $78.96\%$ tail risk reduction at an upfront net debit of only \rupee 2.50/share ($0.88\%$ of portfolio value).
\end{abstract}

\vspace{4pt}
\noindent\textbf{Keywords:} Value at Risk ($\VaR$), Expected Shortfall ($\CVaR$), Kupiec Backtest, Black--Scholes--Merton, Continuous Dividend Yield, Protective Put, Zero-Cost Collar, ITC Limited.

\section{Introduction and the Dual Problem}
\label{sec:introduction}

In institutional portfolio management, market risk quantification and derivative financial engineering are frequently treated as isolated quantitative workflows. Risk managers estimate Value at Risk ($\VaR$) to comply with statutory capital adequacy requirements (such as the Basel II/III market risk frameworks) without translating statistical loss thresholds into concrete derivative hedging positions. Conversely, options traders price derivative contracts using continuous-time arbitrage models without anchoring strike selections to empirical tail loss quantiles.

This paper bridges this institutional disconnect by formulating an integrated, empirical risk architecture applied to an equity position in \textbf{ITC Limited (NSE: ITC)}. Specifically, the investigation resolves the \textit{Dual Quantitative Problem}:
\begin{enumerate}[label=(\roman*), noitemsep]
    \item \textbf{Downside Market Risk Quantification:} Quantifying the maximum potential capital loss over 1-day and 10-day horizons for a standard National Stock Exchange (NSE) F\&O market lot ($1,725$ shares) using Parametric (Delta-Normal) and Non-Parametric Historical Simulation $\VaR$ alongside Expected Shortfall ($\CVaR$), subject to formal Kupiec Proportion of Failures ($\POF$) backtesting.
    \item \textbf{Derivative Hedging Optimization:} Evaluating whether theoretical option pricing models---contrasting Standard Black--Scholes ($q=0$) against Merton's (1973) Continuous Dividend-Adjusted Model ($q=5.09\%$) and Discrete Dividend models---accurately price the derivative instruments required to insulate the shareholder against that quantified tail loss.
\end{enumerate}

\begin{definition}{Downside Tail Risk and Derivative Hedging Duality}{tail_risk_duality}
Let $V_0 = L \cdot S_0$ represent the initial capital commitment of holding $L$ shares at spot price $S_0$. The downside market risk at confidence level $1 - \alpha$ over holding horizon $T_h$ is defined by $\VaR_{\alpha, T_h}$ such that $\Prob(\Delta V \le -\VaR_{\alpha, T_h}) = \alpha$. The derivative hedging objective consists of constructing a non-linear terminal payoff $\Phi(S_T)$ such that the combined portfolio value $\Pi_T = L \cdot S_T + M \cdot \Phi(S_T)$ satisfies a guaranteed minimum capital floor: $\Pi_T \ge V_0 - \mathcal{L}_{\max}, \; \forall S_T \in [0, \infty)$, where $\mathcal{L}_{\max}$ is the maximum acceptable portfolio loss.
\end{definition}
